\documentclass[12pt,a4paper]{ctexart}

\usepackage{amsmath,amssymb,amsthm,mathtools,bm}
\usepackage{geometry}
\usepackage{fancyhdr}
\usepackage{enumitem}
\usepackage[colorlinks=true,linkcolor=blue,citecolor=blue,urlcolor=blue]{hyperref}

\geometry{
  left=2.6cm,
  right=2.6cm,
  top=2.5cm,
  bottom=2.6cm,
  headheight=16pt
}

\setlength{\parindent}{2em}
\setlength{\parskip}{0pt}
\linespread{1.25}

\fancyhf{}
\lhead{山东大学}
\chead{逆映射定理课程报告}
\rhead{泰山学堂}
\cfoot{\thepage}
\renewcommand{\headrulewidth}{0.4pt}
\pagestyle{fancy}

\ctexset{
  section = {
    format = \raggedright\bfseries\zihao{-3},
    beforeskip = 1.2ex,
    afterskip = 0.8ex,
    number = {}
  },
  subsection = {
    format = \raggedright\bfseries\zihao{4},
    beforeskip = 1ex,
    afterskip = 0.6ex,
    number = {}
  }
}

\begin{document}

\section*{摘要}

逆映射定理是多元微分学中最重要的两个定理之一，在本文的前半部分总结整理了李良攀老师的证明，在证明逆函数可微是使用可微性的 Hadamard 刻画定理并结合直接证明可微并得到其微分。在证明可微时我使用利用向量值 Newton-Leibniz 证明局部为单射的证明中时使用的一个推论简化了论证过程。本文的第二部分给出了两个相关命题的加强与拓展，证明过程充分结合了 $\mathbb{R}^n$ 的紧性与完备性与集合联通的性质。

\section*{逆映射定理}

\subsection*{定理1：向量值Newton-Leibniz公式}

\textbf{定理内容：}

设 $U \subset \mathbb{R}^n$ 是一个开集，函数 $f: U \to \mathbb{R}^m$ 是一阶连续可微的（即 $f \in C^1(U, \mathbb{R}^m)$）。若连接 $x$ 与 $x+h$ 的线段 $[x, x+h] = \{x + th \mid t \in [0, 1]\}$ 完全包含在 $U$ 内，则有：

\[
f(x+h) - f(x) = \left( \int_0^1 f'(x+th) \, dt \right) h
\]

其中 $f'$ 表示 $f$ 的 Jacobi 矩阵（即导数），积分是对矩阵的每一个元素分别进行积分。

\textbf{证明：}
构造辅助函数 $\varphi: [0, 1] \to \mathbb{R}^m$，定义为：

\[
\varphi(t) = f(x+th)
\]

由于 $f$ 是连续可微的，并且线段 $[x, x+h] \subset U$，根据多元复合函数的链式法则，$\varphi(t)$ 在 $[0, 1]$ 上连续可微，且其导数为：

\[
\varphi'(t) = f'(x+th) \cdot \frac{d}{dt}(x+th) = f'(x+th)h
\]

这里 $f'(x+th)$ 是一个 $m \times n$ 的矩阵，$h$ 是一个 $n \times 1$ 的列向量，因此 $\varphi'(t)$ 是一个 $m \times 1$ 的列向量。

对向量值函数 $\varphi(t)$ 的每一个分量 $\varphi_i(t)$ ($i = 1, 2, \dots, m$) 分别应用一元函数的 Newton-Leibniz 公式：

\[
\varphi_i(1) - \varphi_i(0) = \int_0^1 \varphi_i'(t) \, dt
\]

将其重新组合为向量形式，有：

\[
\varphi(1) - \varphi(0) = \int_0^1 \varphi'(t) \, dt
\]

注意到 $\varphi(1) = f(x+h)$ 且 $\varphi(0) = f(x)$，代入导数表达式得到：

\[
f(x+h) - f(x) = \int_0^1 f'(x+th)h \, dt
\]

由于 $h$ 与积分变量 $t$ 无关，我们可以利用矩阵乘法和积分的线性性质，将 $h$ 提出积分号外，最终得到：

\[
f(x+h) - f(x) = \left( \int_0^1 f'(x+th) \, dt \right) h
\]

证毕。

\subsection*{定义1：非退化矩阵}
一个 $n \times n$ 的实矩阵 $A$ 称为非退化矩阵（或可逆矩阵），如果$\forall x\in V(\dim V=n):Ax=0 \Rightarrow x=0$。

矩阵$A$非退化与以下命题等价：
\begin{enumerate}[label=\arabic*.,leftmargin=2.5em]
\item 存在唯一的逆矩阵 $A^{-1}$，使得 $A A^{-1} = A^{-1} A = I_n$，其中 $I_n$ 是 $n$ 阶单位矩阵。
\item 它的行列式不为零，即 $\det(A) \neq 0$
\end{enumerate}

在多元微积分中，如果向量值函数 $f: U \to \mathbb{R}^n$ 在点 $x_0$ 处的 Jacobi 矩阵 $f'(x_0)$ 是非退化矩阵，则称 $f$ 在 $x_0$ 处的微分为非退化的。

\subsection*{定理2 ：微分(Jacobi矩阵)某点非退化导出局部为单射}

\textbf{命题内容：}
设 $U \subset \mathbb{R}^n$ 是一个开集，函数 $f: U \to \mathbb{R}^n$ 是一阶连续可微的（即 $f \in C^1(U, \mathbb{R}^n)$）。若 $f$ 在点 $x_0 \in U$ 处的 Jacobi 矩阵 $f'(x_0)$ 是非退化的（即可逆），则存在 $x_0$ 的一个开邻域 $V \subset U$，使得 $f$ 在 $V$ 上是单射。

\textbf{证明：}
由于 $f'(x_0)$ 是非退化矩阵，其逆矩阵 $(f'(x_0))^{-1}$ 存在。设矩阵的算子范数为 $\|\cdot\|$，令 $\lambda = \frac{1}{2 \| (f'(x_0))^{-1} \|} > 0$。

因为 $f \in C^1(U, \mathbb{R}^n)$，Jacobi 矩阵映射 $x \mapsto f'(x)$ 在 $x_0$ 处是连续的。因此，存在 $x_0$ 的一个凸的开球邻域 $V = B(x_0, \delta) \subset U$，使得对任意 $x \in V$，都有：

\[
\| f'(x) - f'(x_0) \| < \lambda
\]

对于任意 $x_1, x_2 \in V$ 且 $x_1 \neq x_2$，由于开球 $V$ 是凸集，连接 $x_1$ 和 $x_2$ 的线段 $[x_1, x_2]$ 完全包含在 $V$ 内。根据定理1（向量值 Newton-Leibniz 公式），我们有：

\[
f(x_2) - f(x_1) = \left( \int_0^1 f'(x_1 + t(x_2 - x_1)) \, dt \right) (x_2 - x_1)
\]

将其改写为：

\[
f(x_2) - f(x_1) = f'(x_0)(x_2 - x_1) + \left( \int_0^1 [f'(x_1 + t(x_2 - x_1)) - f'(x_0)] \, dt \right) (x_2 - x_1)
\]

两边左乘 $(f'(x_0))^{-1}$，得到：

\[
(f'(x_0))^{-1}[f(x_2) - f(x_1)] = (x_2 - x_1) + (f'(x_0))^{-1} \left( \int_0^1 [f'(x_1 + t(x_2 - x_1)) - f'(x_0)] \, dt \right) (x_2 - x_1)
\]

移项并两边取范数，利用三角不等式 $\|a - b\| \ge \|a\| - \|b\|$，有：

\[
\| (f'(x_0))^{-1}[f(x_2) - f(x_1)] \| \ge \| x_2 - x_1 \| - \left\| (f'(x_0))^{-1} \left( \int_0^1 [f'(x_1 + t(x_2 - x_1)) - f'(x_0)] \, dt \right) (x_2 - x_1) \right\|
\]

对于右侧的减数部分，由于 $x_1 + t(x_2 - x_1) \in V$，故 $\| f'(x_1 + t(x_2 - x_1)) - f'(x_0) \| < \lambda$。利用积分范数不等式和矩阵乘法范数的不等式，我们有：

\[
\left\| \int_0^1 [f'(x_1 + t(x_2 - x_1)) - f'(x_0)] \, dt \right\| \le \int_0^1 \| f'(x_1 + t(x_2 - x_1)) - f'(x_0) \| \, dt < \lambda
\]

因此，减数部分的范数满足：

\[
\left\| (f'(x_0))^{-1} \left( \int_0^1 [f'(x_1 + t(x_2 - x_1)) - f'(x_0)] \, dt \right) (x_2 - x_1) \right\| \le \| (f'(x_0))^{-1} \| \cdot \lambda \cdot \| x_2 - x_1 \|
\]

由 $\lambda$ 的定义知 $\| (f'(x_0))^{-1} \| \cdot \lambda = \frac{1}{2}$，所以上式严格小于 $\frac{1}{2} \| x_2 - x_1 \|$（只要 $x_1 \neq x_2$）。

代回原不等式得到：

\[
\| (f'(x_0))^{-1} \| \cdot \| f(x_2) - f(x_1) \| \ge \| (f'(x_0))^{-1}[f(x_2) - f(x_1)] \| > \| x_2 - x_1 \| - \frac{1}{2} \| x_2 - x_1 \| = \frac{1}{2} \| x_2 - x_1 \|
\]

即：

\[
\| f(x_2) - f(x_1) \| > \frac{1}{2 \| (f'(x_0))^{-1} \|} \| x_2 - x_1 \| = \lambda \| x_2 - x_1 \|
\]

因为 $x_1 \neq x_2$，即 $\| x_2 - x_1 \| > 0$，且 $\lambda > 0$，所以必有 $\| f(x_2) - f(x_1) \| > 0$，即 $f(x_1) \neq f(x_2)$。

这就证明了 $f$ 在开邻域 $V$ 上是单射。证毕。

接下来我们来证明一个关键的前置定理：开映射定理

\subsection*{定理3：开映射定理（局部满射性）}

\textbf{命题内容}：
设 $U \subset \mathbb{R}^n$ 是开集，函数 $f: U \to \mathbb{R}^n$ 是 $C^1$ 类映射。若 $f$ 在点 $x_0 \in U$ 处的 Jacobi 矩阵 $J_f(x_0)$ 是非退化的（即非奇异），则 $f(x_0)$ 是像集 $f(U)$ 的内点。也就是说，存在 $x_0$ 的开球 $B(x_0, r) \subset U$ 以及 $f(x_0)$ 的开球 $B(f(x_0), \delta)$，使得 $B(f(x_0), \delta) \subset f(B(x_0, r))$。

\textbf{证明：}
由于 $J_f(x_0)$ 非奇异，且 $f \in C^1(U, \mathbb{R}^n)$，由连续性可知，存在 $x_0$ 的一个闭球 $\overline{B}(x_0, r) \subset U$，使得对所有 $x \in \overline{B}(x_0, r)$，$J_f(x)$ 都是非奇异的，且根据命题1可要求 $f$ 在 $\overline{B}(x_0, r)$ 上是单射。

记 $y_0 = f(x_0)$。由于 $\partial B(x_0, r)$ 是紧集，且 $f$ 在此球面上没有取到 $y_0$（单射性保证 $f(x) \neq y_0, \forall x \in \partial B(x_0, r)$），我们可以定义：
\[
S = \min_{|x-x_0|=r} |f(x) - y_0| > 0
\]

任取 $y \in B(y_0, S/3)$，固定 $y$。我们在闭球 $\overline{B}(x_0, r)$ 上构造关于距离平方的连续函数：
\[
g(x) = |f(x) - y|^2, \quad x \in \overline{B}(x_0, r)
\]

由于 $f(x_0) = y_0$，在中心点 $x_0$ 处有：
\[
g(x_0) = |f(x_0) - y|^2 = |y_0 - y|^2 < \left(\frac{S}{3}\right)^2
\]
即 $|f(x_0) - y| < \frac{S}{3}$.

对于边界上的任意点 $x \in \partial B(x_0, r)$，根据三角不等式：
\[
|f(x) - y| \ge |f(x) - y_0| - |y_0 - y| \ge S - \frac{S}{3} = \frac{2S}{3}
\]
由此可得：
\[
|f(x) - y| \ge \frac{2S}{3} > \frac{S}{3} > |f(x_0) - y|
\]
这表明，对于所有 $x \in \partial B(x_0, r)$，都有 $g(x) > g(x_0)$。

于是，$g(x)$ 的最小值绝对不可能在边界 $\partial B(x_0, r)$ 上取得。
而 $\overline{B}(x_0, r)$ 是有界闭集（紧集），连续函数 $g(x)$ 必在此闭集上取得最小值。
故记 $z_y$ 为 $g(x)$ 取到最小值时 $x$ 的取值，即：
\[
\min_{x \in \overline{B}(x_0, r)} g(x) = g(z_y)
\]
由前面的分析可知，$z_y$ 必定在开球内部，即 $z_y \in B(x_0, r)$。

由于 $z_y$ 是内部极小值点，且 $g(x)$ 可导，由极值的必要条件，其偏导数（梯度）必定为零：
\[
\nabla g(z_y) = 0
\]
根据多元函数链式法则，有：
\[
2 J_f(z_y)^T (f(z_y) - y) = 0
\]

因为 $z_y \in B(x_0, r)$，由我们对 $r$ 的选取，雅可比矩阵 $J_f(z_y)$ 是非奇异的，从而其转置矩阵 $J_f(z_y)^T$ 也是非奇异的。
对上式左乘 $(J_f(z_y)^T)^{-1}$，即可得到：
\[
f(z_y) - y = 0 \implies f(z_y) = y
\]

这证明了对于任意 $y \in B(y_0, S/3)$，都在 $B(x_0, r)$ 内存在原像 $z_y$。
因此：
\[
B(y_0, S/3) \subset f(B(x_0, r)) \subset \text{range}(f)
\]
即 $y_0 = f(x_0)$ 是像集的一个内点，开映射定理得证。

有了局部单射性（命题1）和局部满射性（开映射定理/命题2）作为基础，我们现在给出逆映射定理的完整内容及证明。

\subsection*{定理4：逆映射定理}

\textbf{定理内容：}
设 $U \subset \mathbb{R}^n$ 是开集，函数 $f: U \to \mathbb{R}^n$ 是连续可微的（即 $f \in C^1(U, \mathbb{R}^n)$）。如果 $f$ 在点 $x_0 \in U$ 处的 Jacobi 矩阵 $f'(x_0)$ 是非退化的（即 $\det f'(x_0) \neq 0$），那么：
\begin{enumerate}[label=\arabic*.,leftmargin=2.5em]
\item \textbf{局部双射}：存在 $x_0$ 的开邻域 $V \subset U$ 以及 $y_0 = f(x_0)$ 的开邻域 $W \subset \mathbb{R}^n$，使得限制映射 $f|_V : V \to W$ 是一个双射。
\item \textbf{逆映射连续可微}：其逆映射 $f^{-1}: W \to V$ 也是连续可微的（即 $f^{-1} \in C^1(W, \mathbb{R}^n)$）。
\item \textbf{逆映射的导数公式}：对于任意 $y \in W$（设 $y = f(x)$），逆映射的 Jacobi 矩阵满足：
  \[
  (f^{-1})'(y) = [f'(x)]^{-1}
  \]
\end{enumerate}

\textbf{证明：}

\textbf{第1步：局部双射的建立}
由 $f'(x_0)$ 的连续性和非退化性，根据\textbf{命题1}，存在 $x_0$ 的一个开邻域 $V_1 \subset U$，使得 $f$ 在 $V_1$ 上是单射，并且对所有 $x \in V_1$，$f'(x)$ 均非退化。
根据\textbf{命题2（开映射定理）}，由于 $f$ 在 $V_1$ 上满足条件，像集 $W = f(V_1)$ 是一个开集。
令 $V = V_1$，则 $f: V \to W$ 既是单射又是满射，因此是一个双射。这保证了逆映射 $f^{-1}: W \to V$ 的存在性，且 $W$ 是 $y_0$ 的开邻域。

\textbf{第2步：逆映射的连续性}
设 $w_0 \in W$ 且 $f(z_0) = w_0$，我们需要证明 $f^{-1}$ 在 $w_0$ 处连续。
任取 $\varepsilon > 0$，不妨设 $\varepsilon$ 足够小使得 $B(z_0, \varepsilon) \subset V$。
由\textbf{命题2（开映射定理）}可知，开集 $B(z_0, \varepsilon)$ 的像 $f(B(z_0, \varepsilon))$ 仍然是一个开集。
由于 $w_0 = f(z_0) \in f(B(z_0, \varepsilon))$，作为开集内的点，必然存在一个 $\delta > 0$，使得：
\[
B(w_0, \delta) \subset f(B(z_0, \varepsilon))
\]
由于 $f$ 在 $V$ 上是单射（存在逆映射 $f^{-1}$），对上式两边作用 $f^{-1}$，我们得到：
\[
f^{-1}(B(w_0, \delta)) \subset B(z_0, \varepsilon)
\]
这恰好满足了连续性的 $\varepsilon-\delta$ 定义，即当 $\|w - w_0\| < \delta$ 时，有 $\|f^{-1}(w) - f^{-1}(w_0)\| < \varepsilon$。
这样就证明了逆映射 $f^{-1}$ 在 $W$ 上的连续性。

\textbf{第3步：逆映射的可微性及导数公式}
固定 $y \in W$，设对应唯一的 $x \in V$ 满足 $f(x) = y$。对于 $W$ 中充分靠近 $y$ 的点 $y + k$，存在唯一的 $x + h \in V$ 使得 $f(x + h) = y + k$。
于是有 $k = f(x + h) - f(x)$。
因为 $f$ 在 $x$ 处可微，我们有：
\[
f(x + h) - f(x) = f'(x)h + o(\|h\|) \quad (h \to 0)
\]
即
\[
k = f'(x)h + o(\|h\|)
\]

由于 $f'(x)$ 是可逆的，左乘 $[f'(x)]^{-1}$ 得到：
\[
[f'(x)]^{-1}k = h + [f'(x)]^{-1}o(\|h\|)
\]
移项得：
\[
f^{-1}(y + k) - f^{-1}(y) = h = [f'(x)]^{-1}k - [f'(x)]^{-1}o(\|h\|)
\]

为了证明 $f^{-1}$ 在 $y$ 处可微，我们需要证明当 $k \to 0$ 时，上式尾项是 $o(\|k\|)$。
由\textbf{命题1}的证明过程中我们已知，在 $V$ 上存在常数 $\lambda > 0$ 使得 $\|f(x + h) - f(x)\| \ge \lambda \|h\|$，即 $\|k\| \ge \lambda \|h\|$。
这表明当 $k \to 0$ 时，必定有 $h \to 0$（这也再次印证了连续性），且：
\[
\frac{\|h\|}{\|k\|} \le \frac{1}{\lambda}
\]
是有界的。

于是对于尾项，我们有：
\[
\frac{\|[f'(x)]^{-1}o(\|h\|)\|}{\|k\|} \le \|[f'(x)]^{-1}\| \cdot \frac{o(\|h\|)}{\|h\|} \cdot \frac{\|h\|}{\|k\|} \le \frac{\|[f'(x)]^{-1}\|}{\lambda} \cdot \frac{o(\|h\|)}{\|h\|}
\]
当 $k \to 0$ 时，$h \to 0$，从而 $\frac{o(\|h\|)}{\|h\|} \to 0$。所以尾项确实是 $o(\|k\|)$。

因此，$f^{-1}$ 在 $y$ 处可微，并且其导数（Jacobi 矩阵）为：
\[
(f^{-1})'(y) = [f'(x)]^{-1}
\]

\textbf{第4步：导数的连续性（即 $f^{-1} \in C^1$）}
最后我们需要证明映射 $y \mapsto (f^{-1})'(y)$ 是连续的。
因为 $(f^{-1})'(y) = [f'(f^{-1}(y))]^{-1}$，这一映射可以看作是以下三个映射的复合：
\begin{enumerate}[label=\arabic*.,leftmargin=2.5em]
\item $y \mapsto f^{-1}(y)$ （这在第2步已证连续）
\item $x \mapsto f'(x)$ （这是已知的，因为 $f \in C^1$）
\item 矩阵求逆映射 $A \mapsto A^{-1}$ （对于非退化矩阵，其求逆运算是连续的，因为逆矩阵的元素是原矩阵元素的有理函数且分母不为零）
\end{enumerate}
由于连续映射的复合仍然是连续映射，故 $(f^{-1})'$ 在 $W$ 上连续。

综上所述，逆映射定理得证。

\section*{几个相关命题的加强与拓展}
\subsection*{命题1：(定理3中估计的加强)}

\textbf{命题内容}

假设 $B = B(x_0, r)$ 是 $\mathbb{R}^n$ 中的开球，并且
\[
f : \overline{B} \to \mathbb{R}^n
\]
是一个连续映射，使得 $f$ 在 $B$ 上是 $C^1$ 类、单射且非退化的。

记
\[
y_0 = f(x_0), \qquad S = \min_{|x - x_0|=r} |f(x) - y_0|.
\]

证明：
\[
\overline{B}(y_0, S) \subset \text{range}(f).
\]

在定理三中，我们取得 $\frac{S}{3}$ 并不是特殊的，不难发现，把 $\frac{S}{3}$ 更改为任意更小的实数结论都成立，所以我们考虑将问题转化到一维空间中并使用 $\mathbb{R}$ 的紧性与完备性证明

\textbf{证明}

任取 $y \in \overline{B}(y_0, S)$，我们构造连接 $y_0$ 和 $y$ 的线段：
\[
y_t := y_0 + t(y - y_0), \quad t \in [0, 1]
\]

记 $A := \{t \in [0, 1] : y_t \in f(\overline{B})\}$。我们分四步证明 $A = [0, 1]$

\textbf{第1步: $A$ 非空}

因为 $y_0 = f(x_0)$ 且 $x_0 \in B \subset \overline{B}$，所以 $y_0 \in f(\overline{B})$。当 $t=0$ 时，$y_0 \in f(\overline{B})$，故 $0 \in A$。

\textbf{第2步: $A$ 是 $[0, 1]$ 中的闭集}

对于任意收敛于 $t$ 的点列 $\{t_k\} \subset A$（即 $t_k \to t$），由于 $y_{t_k} \in f(\overline{B})$，我们能够找到原像序列 $x_k \in \overline{B}$ 使得 $f(x_k) = y_{t_k}$。
由于 $\overline{B}$ 是 $\mathbb{R}^n$ 中的紧集（有界闭集），根据 Bolzano–Weierstrass 定理，序列 $\{x_k\}$ 必存在收敛的子列 $\{x_{k_s}\}$，设其收敛于 $x \in \overline{B}$
由 $f$ 的连续性，有：
\[
f(x_{k_s}) \to f(x)
\]
另一方面，由于 $t_{k_s} \to t$，根据线段定义有：
\[
y_{t_{k_s}} \to y_t
\]
因此必有 $f(x) = y_t$。这说明 $y_t \in f(\overline{B})$，从而 $t \in A$。所以 $A$ 是闭集。

\textbf{第3步: $A$ 在 $[0, 1)$ 中是开集}

对任意 $t \in A \cap [0, 1)$，存在 $x_t \in \overline{B}$ 使得 $f(x_t) = y_t$。
我们估计 $f(x_t)$ 到 $y_0$ 的距离：
\[
|f(x_t) - y_0| = |y_t - y_0| = |t(y - y_0)| = t|y - y_0| \le tS < S
\]
（这里利用了 $y \in \overline{B}(y_0, S) \Rightarrow |y - y_0| \le S$ 且 $t < 1$）
由 $S$ 的定义，$S = \min_{|x - x_0|=r} |f(x) - y_0|$，即边界 $\partial B$ 上的点映射后到 $y_0$ 的距离至少为 $S$。因为 $|f(x_t) - y_0| < S$，必定有 $x_t \notin \partial B$。
由于 $x_t \in \overline{B}$ 且 $x_t \notin \partial B$，故 $x_t \in B$。
因为 $f$ 在开球 $B$ 上是 $C^1$ 类且非退化的，根据开映射定理，像集 $f(B)$ 在 $f(x_t) = y_t$ 的某邻域内是开的。
即存在充分小的 $\delta > 0$，使得：
\[
B(y_t, \delta) \subset f(B) \subset f(\overline{B})
\]
由此可知，存在包含 $t$ 的一个小区间 $(t - \varepsilon, t + \varepsilon)$，使得对应的 $y_{t'}$ 均落入 $B(y_t, \delta)$ 内，从而 $y_{t'} \in f(\overline{B})$。这证明了 $A$ 在 $[0, 1)$ 中是开集

\textbf{第4步: 结论导出}

令 $a := \sup A$。由于 $A$ 在 $[0, 1]$ 中闭，必定有 $a \in A$。
假设 $a < 1$，由 Step 3 知 $A$ 在 $[0, 1)$ 中是开集，既然 $a \in [0, 1) \cap A$，那么 $A$ 必然包含 $a$ 的一个右侧小区间 $(a, a + \varepsilon)$。但这与 $a$ 是上确界矛盾。
因此必然有 $\sup A = 1$
又因为 $A$ 闭，所以 $1 \in A$。
当 $t = 1$ 时，$y_1 = y_0 + (y - y_0) = y$。因此 $y \in f(\overline{B})$。
由于 $y \in \overline{B}(y_0, S)$ 是任取的，我们最终得到：
\[
\overline{B}(y_0, S) \subset f(\overline{B}) = \text{range}(f)
\]


\subsection*{命题2：考察映射f有孤立奇异点的情况下能否保证开映射}

\textbf{问题描述：}
假设 $B(x_0, \delta)$ 是 $\mathbb{R}^2$ 中的开球，$f: B(x_0, \delta) \to \mathbb{R}^2$ 是一个连续映射。
记 $\Omega^* = B(x_0, \delta) \setminus \{x_0\}$。
如果限制在 $\Omega^*$ 上的映射 $f|_{\Omega^*}$ 是一个没有临界点（without CPs，即雅可比行列式处处不为零的 $C^1$ 映射），那么 $f(B(x_0, \delta))$ 在 $\mathbb{R}^2$ 中是否是开集？

我们考虑在顺着命题2的思路给出这个命题的证明
不失一般性，通过平移我们可以假设 $x_0 = 0$ 且 $f(0) = 0$。我们需要证明像点 $0$ 是像集 $f(B(0, \delta))$ 的内点

\textbf{第1步：像集在去心邻域是开集}

由于 $f$ 在去心邻域 $\Omega^* = B(0, \delta) \setminus \{0\}$ 上是 $C^1$ 且雅可比行列式非零，根据逆映射定理（或开映射定理），$f$ 在 $\Omega^*$ 上是开映射。因此，$f(\Omega^*)$ 是 $\mathbb{R}^2$ 中的开集。

\textbf{第2步：转化为内点问题}

要证明 $f(B(0, \delta)) = f(\Omega^*) \cup \{0\}$ 是开集，因为 $f(\Omega^*)$ 已知是开集，我们只需证明 $0$ 是 $f(B(0, \delta))$ 的内点。即证明存在 $S > 0$，使得 $\overline{B}(0, S) \subset f(B(0, \delta))$。

\textbf{第3步：构造像不含0的边界}

考虑被映成 0 的原像，设 $x^*\in B(0,\delta)\setminus\{0\}$ 满足 $f(x^*)=0$，
由于 $f$ 在 $\Omega^*$ 上是微分同胚，因而局部单射，那么在 $x^*$ 附近可以去一个充分小的开球 $B(x^*,\delta^*)$，使得在这个开球上只有 $x^*$ 一个点被映成 0
又由于有理数的稠密性，在这个开球中至少可以找到一个点它的横纵坐标都是有理数；而 $\mathbb{Q}$是可数的，$\mathbb{Q}^2$ 依然可数，这说明被映成0的原像至多可数，包含这些点的圆 $\partial B(0,r)$ 至多可数。但是 $(0,\delta)$ 不可数，那么总存在一些 $r\in(0,\delta)$，使得 $\forall x \in \partial B(0,r):f(x)\not =0$

由于 $\partial B(0, r)$ 是紧集且 $f$ 连续，像集 $f(\partial B(0, r))$ 也是紧集。
定义：
\[
S = \min_{|x|=r} |f(x) - 0| = \min_{|x|=r} |f(x)| > 0
\]

\textbf{第4步：利用序列极限刻画边界}

令 $V_x = B(0, r) \setminus \{0\}$。这是一个开集，根据第 1 步，像集 $U_y = f(V_x)$ 是 $\mathbb{R}^2$ 中的开集。
现在，我们用序列极限来分析这个开集 $U_y$ 的边界 $\partial U_y$

任取边界上的一点 $y \in \partial U_y$。
根据边界的定义，必然存在一个序列 $y_n \in U_y$，使得 $\lim_{n \to \infty} y_n = y$
因为 $y_n \in U_y = f(V_x)$，必定存在原像序列 $x_n \in V_x \subset \overline{B}(0, r)$，使得 $f(x_n) = y_n$。

由于 $x_n$ 是有界闭集 $\overline{B}(0, r)$ 中的序列，根据Bolzano-Weierstrass 定理，它必定存在一个收敛子列 $x_{n_k}$，收敛到某个极限点 $x^* \in \overline{B}(0, r)$。
由 $f$ 的连续性，两边取极限：
\[
f(x^*) = \lim_{k \to \infty} f(x_{n_k}) = \lim_{k \to \infty} y_{n_k} = y
\]

此时我们来判断 $x^*$ 的位置：
\begin{enumerate}[label=\arabic*.,leftmargin=2.5em]
\item 如果 $x^* \in V_x$，由于 $V_x$ 是开集且 $f$ 在 $V_x$ 上是开映射，那么 $y = f(x^*)$ 必须是 $U_y$ 的内点。但这与 $y \in \partial U_y$ 矛盾！
\item 因此，$x^*$ 绝对不能在 $V_x$ 内。
\item 既然 $x^* \in \overline{B}(0, r)$ 且 $x^* \notin V_x$，那么 $x^*$ 只能在 $V_x$ 的边界上。而 $V_x$ 的边界正是圆周 $\partial B(0, r)$ 以及原点 $\{0\}$。
\end{enumerate}

由此，我们用得到了一个关键的边界包含关系：
\[
\partial U_y \subset f(\partial B(0, r)) \cup \{f(0)\} = f(\partial B(0, r)) \cup \{0\}
\]
这说明，像集 $U_y$ 的边界只能来源于原定义域的边界或原点。

\textbf{第5步：利用连通性确定内点}

现在，考察原点周围的开球 $B(0, \frac{S}{2})$。
根据 $S$ 的定义，对于任何 $x \in \partial B(0, r)$，都有 $|f(x)| \ge S$。这意味着 $f(\partial B(0, r))$ 与开球 $B(0, \frac{S}{2})$ 完全没有交集。

结合第 4 步的结论，在开球 $B(0, \frac{S}{2})$ 内部，$U_y$ 的边界点只可能有一个，那就是原点 $0$ 本身。
换言之，去心开球 $W = B(0, \frac{S}{2}) \setminus \{0\}$ 中，绝对不包含 $U_y$ 的任何边界点，即 $W\cap \partial U_y=\emptyset$

在 $\mathbb{R}^2$ 中，去心开球 $W = B(0, \frac{S}{2}) \setminus \{0\}$ 是一个连通集。
既然连通集 $W$ 完全不包含开集 $U_y$ 的边界，那么 $W$ 要么完全包含在 $U_y$ 内，要么完全在 $U_y$ 外。否则，我们可以在边界两边各取一个点 $x_{in},x_{out}$，由界值定理知其连线必然穿过 $U_y$ 的边界，这与 $W\cap \partial U_y=\emptyset$ 矛盾

因为 $f$ 连续且 $f(0) = 0$，当 $x \to 0$ 时，$f(x) \to 0$。  
对于给定的 $\frac{S}{2}$，$\exists \eta \in (0,r),\forall x \in B(0,\eta):f(x) < \frac{S}{2}$
并且 $f(B(0,\eta))$ 不恒为零，否则与 $f$ 只有孤立临界点矛盾。  
这意味着 $U_y = f(V_x)$ 中必然有点落在 $W$ 内，所以 $W$ 不可能完全在 $U_y$ 外。
因此，$W$ 必须完全在 $U_y$ 内部，即：
\[
B(0, S) \setminus \{0\} \subset U_y = f(B(0, r) \setminus \{0\})
\]

最后，加上原点的映射关系 $f(0) = 0$，我们得到：
\[
B(0, S) \subset f(B(0, r)) \subset f(B(0, \delta))
\]
即原点 0 是像集的内点，命题得证。 

\section*{参考文献}
[1] 李良攀. InverseFT. 2024

[2] V. A. Zorich. Mathematical Analysis I.  2016.

\end{document}
